% GENERATED FILE. Edit canonical lessons, then run npm run generate:books.
% canonical-source-hash: fnv1a64:1dd50f341814fca4
% canonical-lessons: FR-C124-boulanger, FR-C124-etudiant, FR-C124-chauffeur, FR-C124-avocat, FR-C124-ingenieur

\chapter{Baker, Student, Driver, Lawyer, Engineer}
\label{ch:fr-baker-student-driver-lawyer-engineer}
\hlchaptermodality{124}

\begin{chapteropening}
\textbf{By the end of this chapter:} \emph{I can name a baker, a student, a driver, a lawyer, an engineer.}

Complete the last lesson of chapter 124: I can name a baker, a student, a driver, a lawyer, an engineer.
\end{chapteropening}

\section[le boulanger]{le boulanger — a baker}

\label{lesson:FR-C124-boulanger}



\begin{quote}
Before the new one: say the French for a shop assistant, then the French for a farmer.
\end{quote}

\subsection*{You'll want to know: le boulanger}
\textbf{le boulanger} — \textquotedblleft{}a baker\textquotedblright{}. The person behind \textbf{la boulangerie}.

\begin{cousinweb}
\textbf{boulanger} $\leftarrow$ Picard \emph{boulenc}, probably from Middle Dutch \emph{bolle}, \textquotedblleft{}a round loaf\textquotedblright{}: a maker of round loaves. (EtymDB 2.1 has its Old French form.)
\end{cousinweb}

\begin{grammarlens}[title={What you've built}]
One more word for this chapter.
\end{grammarlens}

\subsection*{Your turn}
\begin{itemize}
\raggedright
  \item \emph{Say it:} \emph{le boulanger}
  \item \emph{Say it:} \emph{le boulanger}, once more
  \item \emph{Say it:} say \emph{le boulanger}
  \item \emph{Recall:} say the French for a shop assistant, then the French for a farmer, then say \emph{le boulanger} again
\end{itemize}

\begin{tcolorbox}[breakable,colback=teal!4,colframe=teal!35!black,title={Before you move on}]
What does le boulanger mean? (\textquotedblleft{}A baker\textquotedblright{}.) Say it once more.
\end{tcolorbox}
\section[l'étudiant]{l'étudiant — a student}

\label{lesson:FR-C124-etudiant}



\begin{quote}
Before the new one: say the French for a farmer, then the French for a baker.
\end{quote}

\subsection*{You'll want to know: l'étudiant}
\textbf{l'étudiant} — \textquotedblleft{}a student\textquotedblright{}.

\begin{cousinweb}
\textbf{étudiant} is the present participle of \emph{étudier}, \textquotedblleft{}to study\textquotedblright{}, used as a noun: one who is studying. \emph{Étudier} $\leftarrow$ Old French \emph{estudier}, from the Latin family of \emph{studium}, \textquotedblleft{}eagerness, study\textquotedblright{} (EtymDB 2.1).
\end{cousinweb}

\begin{grammarlens}[title={What you've built}]
One more word for this chapter.
\end{grammarlens}

\subsection*{Your turn}
\begin{itemize}
\raggedright
  \item \emph{Say it:} \emph{l'étudiant}
  \item \emph{Say it:} \emph{l'étudiant}, once more
  \item \emph{Say it:} say \emph{l'étudiant}
  \item \emph{Recall:} say the French for a farmer, then the French for a baker, then say \emph{l'étudiant} again
\end{itemize}

\begin{tcolorbox}[breakable,colback=teal!4,colframe=teal!35!black,title={Before you move on}]
What does l'étudiant mean? (\textquotedblleft{}A student\textquotedblright{}.) Say it once more.
\end{tcolorbox}
\section[le chauffeur]{le chauffeur — a driver}

\label{lesson:FR-C124-chauffeur}



\begin{quote}
Before the new one: say the French for a baker, then the French for a student.
\end{quote}

\subsection*{You'll want to know: le chauffeur}
\textbf{le chauffeur} — \textquotedblleft{}a driver\textquotedblright{}.

\begin{cousinweb}
\textbf{chauffeur} — built from the verb for \textquotedblleft{}to heat\textquotedblright{}, with \emph{-eur}: first the stoker of a steam engine, then the driver. (EtymDB 2.1)
\end{cousinweb}

\begin{grammarlens}[title={What you've built}]
One more word for this chapter.
\end{grammarlens}

\subsection*{Your turn}
\begin{itemize}
\raggedright
  \item \emph{Say it:} \emph{le chauffeur}
  \item \emph{Say it:} \emph{le chauffeur}, once more
  \item \emph{Say it:} say \emph{le chauffeur}
  \item \emph{Recall:} say the French for a baker, then the French for a student, then say \emph{le chauffeur} again
\end{itemize}

\begin{tcolorbox}[breakable,colback=teal!4,colframe=teal!35!black,title={Before you move on}]
What does le chauffeur mean? (\textquotedblleft{}A driver\textquotedblright{}.) Say it once more.
\end{tcolorbox}
\section[l'avocat]{l'avocat — a lawyer; an avocado}

\label{lesson:FR-C124-avocat}



\begin{quote}
Before the new one: say the French for a student, then the French for a driver.
\end{quote}

\subsection*{You'll want to know: l'avocat}
\textbf{l'avocat} — \textquotedblleft{}a lawyer; an avocado\textquotedblright{}.

\begin{cousinweb}
\textbf{avocat} — two words in one. The lawyer $\leftarrow$ Latin \emph{advocatus}, \textquotedblleft{}one called in to help\textquotedblright{} (EtymDB 2.1) — a learned borrowing. The fruit $\leftarrow$ Spanish \emph{aguacate}, from Nahuatl \emph{ahuacatl} (EtymDB 2.1), reshaped to look like the lawyer word.
\end{cousinweb}

\begin{grammarlens}[title={What you've built}]
One more word for this chapter.
\end{grammarlens}

\subsection*{Your turn}
\begin{itemize}
\raggedright
  \item \emph{Say it:} \emph{l'avocat}
  \item \emph{Say it:} \emph{l'avocat}, once more
  \item \emph{Say it:} say \emph{l'avocat}
  \item \emph{Recall:} say the French for a student, then the French for a driver, then say \emph{l'avocat} again
\end{itemize}

\begin{tcolorbox}[breakable,colback=teal!4,colframe=teal!35!black,title={Before you move on}]
What does l'avocat mean? (\textquotedblleft{}A lawyer; an avocado\textquotedblright{}.) Say it once more.
\end{tcolorbox}
\section[l'ingénieur]{l'ingénieur — an engineer}

\label{lesson:FR-C124-ingenieur}



\begin{quote}
Before the new one: say the French for a driver, then the French for a lawyer; an avocado.
\end{quote}

\subsection*{You'll want to know: l'ingénieur}
\textbf{l'ingénieur} — \textquotedblleft{}an engineer\textquotedblright{}.

\begin{cousinweb}
\textbf{ingénieur} — built from Old French \emph{engin}, \textquotedblleft{}a clever device, a war machine\textquotedblright{}, $\leftarrow$ Latin \emph{ingenium}, \textquotedblleft{}talent\textquotedblright{} (EtymDB 2.1), with \emph{-eur}, later respelt after the Latin.
\end{cousinweb}

\begin{grammarlens}[title={What you've built}]
One more word for this chapter.
\end{grammarlens}

\subsection*{Your turn}
\begin{itemize}
\raggedright
  \item \emph{Say it:} \emph{l'ingénieur}
  \item \emph{Say it:} \emph{l'ingénieur}, once more
  \item \emph{Say it:} say \emph{l'ingénieur}
  \item \emph{Recall:} say the French for a driver, then the French for a lawyer; an avocado, then say \emph{l'ingénieur} again
\end{itemize}

\begin{tcolorbox}[breakable,colback=teal!4,colframe=teal!35!black,title={Before you move on}]
What does l'ingénieur mean? (\textquotedblleft{}An engineer\textquotedblright{}.) Say it once more.
\end{tcolorbox}
